% ID: SEL-D0028-OA-E01
% ESTADO: cotejada. Segunda pasada visual completa y borrador compilado revisados.
% Fuente: Selectividad/matematicas_aplicadas/suplente junio/Examen-Matemáticas Aplicadas a las Ciencias Sociales.pdf
% SHA-256: a10146fce64a17e21450973c88d66d43571619781f76bc345b5a4b3dee67a16b
% Página: 1; opción A; ejercicio original 1.

Se considera la región definida por las siguientes inecuaciones:
\[
2x-y\geq2\qquad -x+2y\leq2\qquad3x+y\leq15\qquad y\geq0
\]
\begin{enumerate}[label=\alph*)]
\item \textbf{(1.8 puntos)} Represéntela gráficamente y determine sus vértices.
\item \textbf{(0.2 puntos)} Indique razonadamente si el punto $(3,3)$ pertenece a dicha región.
\item \textbf{(0.5 puntos)} \textquestiondown{}En qué puntos de la región anterior la función $F(x,y)=3x-2y$ alcanza los valores máximo y mínimo y cuáles son éstos?
\end{enumerate}
